% Preview of the MIPS-algorithms appendix text + combined figure.
% Build: python figures/make_mips_algorithms_figure.py && cd figures && tectonic mips_algorithms_appendix.tex
% Page geometry approximates an ICLR page (5.5in x 9in text block).
\documentclass[10pt]{article}
\usepackage[textwidth=5.5in,textheight=9in]{geometry}
\usepackage{amsmath,amssymb,booktabs,xcolor,tabularx,parskip}
\usepackage[hidelinks]{hyperref}
\newcommand{\op}[1]{\mathrm{#1}}
\begin{document}

%% ===================== BEGIN APPENDIX TEXT (paste from here) =====================
\section{Decoding the recovered MIPS programs}\label{app:mips-algorithms}

Figure~\ref{fig:mips-algorithms} shows the programs recovered by MIPS-evolved PySR on
Parity-last~4, Alternating-last~4, and Base-5 addition. Throughout, $h=(h_0,h_1,\ldots)$ is
the RNN hidden state after projection onto its learned integer lattice, primes denote the
next state, $x$ (or the digit pair $a,b$) is the current input, and $y$ is the output read
from the new state. To interpret each program we labeled every lattice state with the input
histories that reach it, and then checked each recovered equation case by case against that
labeling.

\paragraph{Parity-last 4.} The algorithm returns the parity of the last four bits of the input.
The hidden state has three coordinates, $h_0\in\{0,2,3,4,8\}$, $h_1\in\{0,1,2,5,7,9\}$, and
$h_2\in\{0,4\}$; the input $x$ and output $y$ are binary. Of all combinations of these values,
only 16 states are observed, and they correspond exactly to the 16 windows $d\,a\,b\,c$ of the
last four input bits (oldest to newest). The pair $(h_0,h_1)$ is a shift register over the
last three bits $a\,b\,c$:
\[
\begin{array}{c|cccccccc}
abc & 000 & 100 & 010 & 110 & 001 & 101 & 011 & 111\\\hline
(h_0,h_1) & (2|3,0) & (0,1) & (0,2) & (4,2) & (0,5) & (3|4,5) & (8,7) & (4,9)
\end{array}
\]
where $2|3$ and $3|4$ denote two lattice points for the same three bits, depending on the bit
$d$ that is about to be dropped. The coordinate $h_2$ stores the parity of all four bits.

\emph{Parity update.} When $x$ arrives, the new window is $a\,b\,c\,x$, so the new parity is
$\op{parity}(a,b,c)\oplus x$. With the embedding of $(a, b, c)$ into $(h_0, h_1)$, the three-bit parity is linearly separable:
$\op{parity}(a,b,c)=0$ exactly when $h_1<2h_0$. The recovered update
$h_2'=4\,\op{eq}(\op{lt}(h_1,2h_0),x)$ therefore equals $4\,(\op{parity}(a,b,c)\oplus x)$: the
indicator $\op{lt}(h_1,2h_0)=1-\op{parity}(a,b,c)$ is compared with $x$, and the result is
multiplied by four. The output $y=\min(1,h_2)$ maps $\{0,4\}$ back to $\{0,1\}$. The extra factor of four is an artifact of the lattice refinement during the state space discretization step of the RNN interpretation pipeline.

\emph{Shift-register update.} $h_1$ encodes the last two bits $(b,c)$:
$00\mapsto\{0,1\}$, $10\mapsto2$, $01\mapsto5$, $11\mapsto\{7,9\}$, where the alternatives depend on
the older bit $a$. The recovered $h_1'$ shifts this code and pushes in $x$. The new last two bits
are $(c,x)$, and $c$ can be read off from $h_1$ ($c=1$ iff $h_1\ge5$). For $x=0$ the update reduces
to $\lfloor\min(0.66h_1,2.32)\rfloor$, which maps $\{0,1\}\mapsto0$ ($c\,x=00$), $2\mapsto1$ ($c\,x=00$, with
older bit $1$), and $\{5,7,9\}\mapsto2$ ($c\,x=10$). For $x=1$ it reduces to
$\lfloor\min(2.32+\max(3,h_1),9.32)\rfloor$, which maps $h_1\le2\mapsto5$ ($01$), $5\mapsto7$, and
$\{7,9\}\mapsto9$ (both $11$). The $2.98567x$ term in the recovered expression is inactive, because
the inner minimum saturates at $2.32$ whenever $x=1$. The values of $h_0$ follow no
simple bit pattern, and the recovered $h_0'$ acts as a lookup table that reproduces them from
$(h_1,x)$, using $h_0$ only in the single case where $h_1$ alone is ambiguous.

\paragraph{Alternating-last 4.}
Here $y=1$ iff the last four bits alternate (0101 or 1010). The RNN stores the last bit and
the length $L$ of the current alternating run in $h_0$, and a flag $h_1=\mathbf 1[L\ge4]$, so
$y=h_1$:
\[
\begin{array}{c|ccc|cccc}
h_0 & 6,7,8 & 1 & 0 & 31,32 & 37\,(h_1{=}0) & 38 & 37\,(h_1{=}1)\\\hline
\text{last bit} & 1&1&1 & 0&0&0&0\\
L & 1 & 2 & \ge3 & 1 & 2 & 3 & \ge4
\end{array}
\]
Small values of $h_0$ encode a last bit of 1 and large values (31--38) a last bit of 0; the
states $\{6,7,8\}$ and $\{31,32\}$ are redundant copies with identical transitions.
Up to a $\pm1$ correction, the recovered $h_0'$ is the affine map
$h_0'\approx\lfloor7.83-h_0/5+30\cdot\mathbf 1[x{=}0]\rfloor$, which moves the state between the
two bands. For $x=1$, a large $h_0$ (the run continues) maps to $1$ ($L=2$) or $0$ ($L\ge3$),
and a small $h_0$ (a repeated bit, which resets the run) maps to $\{6,7,8\}$. For $x=0$, the
offset of 30 lands in the large band: a repeated 0 maps to $\{31,32\}$, and a run ending in 1
maps to 37 or 38. The term $\op{eq}(\op{not}(h_0),x)$ corrects the case $h_0=0$, $x=0$, which
reaches $L\ge4$. The recovered $h_1'$ sets the flag exactly when the run already had length
at least three and $x$ continues it: the factor $\op{zero}(\op{xor}(\op{lt}(x,h_0),x))$ is one
iff ($x=0$ and $h_0=0$) or ($x=1$ and $h_0\ge2$), and the second factor keeps only the
$L\ge3$ states ($h_0=0$, $h_0=38$, or $h_0=37$ with $h_1=1$). The old flag $h_1$ is needed
because $h_0=37$ alone does not distinguish $L=2$ from $L\ge4$.

\paragraph{Base-5 addition.}
Each step reads one digit of each summand, $a,b\in\{0,\ldots,4\}$, least significant first,
and outputs the corresponding digit of the sum. The ten states correspond to the ten values
of the column sum $s=a+b+c_{\mathrm{in}}\in\{0,\ldots,9\}$, encoded as a carry bit
$h_1=\mathbf 1[s\ge5]$ and a digit coordinate $h_0=6-s+3h_1$ that decreases in $s$ and jumps
at the carry ($s=0,\ldots,4\mapsto6,\ldots,2$; $s=5,\ldots,9\mapsto4,\ldots,0$). The output
$y=|h_0+h_1-6|-h_1$ inverts this encoding: without a carry it gives $|{-s}|=s$, and with a
carry $|4-s|-1=s-5$, i.e.\ $y=s\bmod5$. The carry update $h_1'=\mathbf 1[a+b+h_1\ge5]$ is
implemented by two branches. Without an incoming carry the mod term vanishes and
$\op{eq}(1,0.119(a+b))$ fires iff $a+b\ge5$. With an incoming carry, the term
$\op{mod}(-1.033,0.254(a+b))$ equals $\approx1.0$ for $a+b=4$ and $0.49$, just below the
$0.5$ threshold of $\op{eq}$, for $a+b\le3$; this is correct on all inputs but relies on a
numerical coincidence. The digit update computes $6-s'+3\cdot\mathbf 1[s'\ge5]$ for
$s'=a+b+h_1$ as the maximum of a linear branch $6.06-s'$, which is correct for $s'<5$, and a
mod branch: for $s'\ge5$, two periods of $4.38$ give $\op{mod}(0.24-s',4.38)\approx9-s'$, the
wrapped code. The same pattern, a linear branch plus a mod branch that implements the
``subtract the base'' wrap-around, appears in the recovered digit updates for base-4 and
base-7 addition.

\paragraph{Summary.}
In all three tasks the RNN representation is a genuine algorithm: a shift register with a
running parity, the last bit with an alternating run length, and the column sum stored as a
carry and a reversed digit. The outputs, the parity rule $h_1<2h_0$, and the carry thresholds
are compact and interpretable. The remaining coordinates are recovered as correct piecewise
lookups built from floor, mod, and max, which often exploit numerical coincidences or the
RNN's arbitrary and partly redundant lattice coordinates; these are precisely the
subproblems that baseline and SRBench-evolved PySR tend to miss.
%% ===================== END APPENDIX TEXT =====================

% Generated by figures/make_mips_algorithms_figure.py -- do not edit by hand.
% Preamble: \usepackage{amsmath,amssymb,booktabs,xcolor,tabularx}
\begin{figure*}[p]
\centering
\begingroup
\definecolor{parityBase}{HTML}{317DA5}
\definecolor{paritySRBench}{HTML}{BB7A20}
\definecolor{parityMIPS}{HTML}{C33D3D}
\providecommand{\op}[1]{\mathrm{#1}}
\footnotesize
\setlength{\tabcolsep}{3pt}
\renewcommand{\arraystretch}{1.35}
\begin{tabularx}{\textwidth}{@{}c ccc X@{}}
\toprule
 & \textcolor{parityBase}{\textbf{Base}} & \textcolor{paritySRBench}{\textbf{Evol.}} & \textcolor{parityMIPS}{\textbf{MIPS}} & \textbf{Discovered equation (MIPS-evolved)} \\
\midrule
\multicolumn{5}{@{}p{\textwidth}@{}}{\textbf{Parity-last 4.} $y=$ parity of the last 4 input bits. $h_0\in\{0,2,3,4,8\}$, $h_1\in\{0,1,2,5,7,9\}$, $h_2\in\{0,4\}$ (16 states), $x,y\in\{0,1\}$. State: $(h_0,h_1)$ holds the last 3 bits, $h_2=4\cdot$parity of the last 4.} \\
\cmidrule(l{0pt}r{0pt}){1-5}
$h_0'$ & \textcolor{gray}{0/10} & \textcolor{gray}{0/10} & \textcolor{parityMIPS}{\textbf{1/10}} & $\begin{aligned}[t]&\min\bigl(((\max(h_1,h_0)-x)h_1-x)h_1-x+2,\\&\;\;4\,|\op{xor}(h_1(0.227249-x),h_1)+x-\op{zero}(h_1)|\bigr)\end{aligned}$ \\
$h_1'$ & \textcolor{parityBase}{\textbf{5/10}} & \textcolor{paritySRBench}{\textbf{3/10}} & \textcolor{parityMIPS}{\textbf{10/10}} & $\begin{aligned}[t]&\bigl\lfloor\min\bigl(\min(2.98567x+0.660363h_1,2.32073)\\&\;\;+x\max(3.00001,h_1),\,9.32075\bigr)\bigr\rfloor\end{aligned}$ \\
$h_2'$ & \textcolor{parityBase}{\textbf{10/10}} & \textcolor{paritySRBench}{\textbf{10/10}} & \textcolor{parityMIPS}{\textbf{10/10}} & $4\,\op{eq}(\op{lt}(h_1,2h_0),x)$ \\
$y$ & \textcolor{parityBase}{\textbf{10/10}} & \textcolor{paritySRBench}{\textbf{10/10}} & \textcolor{parityMIPS}{\textbf{10/10}} & $\min(1,h_2)$ \\
\addlinespace[1pt]
\textit{joint} & \textcolor{gray}{0/10} & \textcolor{gray}{0/10} & \textcolor{parityMIPS}{\textbf{1/10}} & \textit{seeds recovering all subproblems together} \\
\midrule[\heavyrulewidth]
\multicolumn{5}{@{}p{\textwidth}@{}}{\textbf{Alternating-last 4.} $y=1$ iff the last 4 bits alternate. $h_0\in\{0,1,6,7,8,31,32,37,38\}$, $h_1\in\{0,1\}$ (11 states), $x,y\in\{0,1\}$. State: $h_0$ holds the last bit and alternating run length $L$; $h_1=\mathbf 1[L\ge4]$.} \\
\cmidrule(l{0pt}r{0pt}){1-5}
$h_0'$ & \textcolor{parityBase}{\textbf{1/10}} & \textcolor{gray}{0/10} & \textcolor{parityMIPS}{\textbf{5/10}} & $\begin{aligned}[t]&\lfloor\op{eq}(\op{not}(h_0),x)+30\,\op{not}(x)\\&\;\;-0.203372\,h_0+7.82772\rfloor\end{aligned}$ \\
$h_1'$ & \textcolor{parityBase}{\textbf{1/10}} & \textcolor{paritySRBench}{\textbf{10/10}} & \textcolor{parityMIPS}{\textbf{10/10}} & $\begin{aligned}[t]&\op{zero}(\op{xor}(\op{lt}(x,h_0),x))\\&\;\;\cdot\max(h_0-37,\max(\op{zero}(x),h_1))\end{aligned}$ \\
$y$ & \textcolor{parityBase}{\textbf{10/10}} & \textcolor{paritySRBench}{\textbf{10/10}} & \textcolor{parityMIPS}{\textbf{10/10}} & $h_1$ \\
\addlinespace[1pt]
\textit{joint} & \textcolor{gray}{0/10} & \textcolor{gray}{0/10} & \textcolor{parityMIPS}{\textbf{5/10}} & \textit{seeds recovering all subproblems together} \\
\midrule[\heavyrulewidth]
\multicolumn{5}{@{}p{\textwidth}@{}}{\textbf{Base-5 addition.} Inputs $a,b\in\{0,\ldots,4\}$ are digits of two numbers (least significant first); $y\in\{0,\ldots,4\}$ is the sum digit. $h_0\in\{0,\ldots,6\}$, $h_1\in\{0,1\}$ (10 states). State: column sum $s=a+b+c_{\mathrm{in}}$ as carry $h_1=\mathbf 1[s\ge5]$ and digit coordinate $h_0=6-s+3h_1$.} \\
\cmidrule(l{0pt}r{0pt}){1-5}
$h_0'$ & \textcolor{gray}{0/10} & \textcolor{paritySRBench}{\textbf{2/10}} & \textcolor{parityMIPS}{\textbf{3/10}} & $\begin{aligned}[t]&\lfloor\max(6.06243-(a{+}b{+}h_1),\\&\;\;\op{mod}(0.243552-(a{+}b{+}h_1),\,4.38026))\rfloor\end{aligned}$ \\
$h_1'$ & \textcolor{parityBase}{\textbf{8/10}} & \textcolor{paritySRBench}{\textbf{10/10}} & \textcolor{parityMIPS}{\textbf{10/10}} & $\begin{aligned}[t]&\op{eq}(1,\max(0.119279(a{+}b),\\&\;\;\op{mod}(-1.03325,\,0.25413(a{+}b)h_1)))\end{aligned}$ \\
$y$ & \textcolor{parityBase}{\textbf{10/10}} & \textcolor{paritySRBench}{\textbf{10/10}} & \textcolor{parityMIPS}{\textbf{9/10}} & $|h_0+h_1-6|-h_1$ \\
\addlinespace[1pt]
\textit{joint} & \textcolor{gray}{0/10} & \textcolor{paritySRBench}{\textbf{2/10}} & \textcolor{parityMIPS}{\textbf{3/10}} & \textit{seeds recovering all subproblems together} \\
\bottomrule
\end{tabularx}
\endgroup
\caption{Programs recovered from RNNs by MIPS-evolved PySR on three MIPS tasks. Each task is split into scalar regression problems, one per next-state coordinate ($h_i'$) and one for the output ($y$). Counts give exact recovery out of ten search seeds for baseline PySR (Base), SRBench-evolved PySR (Evol.), and MIPS-evolved PySR (MIPS). Equations are regrouped, simplified using the observed value sets, and rounded to six significant digits; each displayed form is verified on the complete recorded relation. Operators: $\op{lt}(u,v)=\mathbf 1[u<v]$, $\op{eq}(u,v)=\mathbf 1[|u-v|<0.5]$, $\op{zero}(u)=\op{not}(u)=\mathbf 1[|u|<0.5]$, $\op{mod}(u,v)=u\bmod|v|$, $\op{xor}(u,v)=(\op{round}(u)+\op{round}(v))\bmod2$. See Appendix~\ref{app:mips-algorithms} for a full decoding of each program.}
\label{fig:mips-algorithms}
\end{figure*}


\end{document}
